\documentclass[12pt]{article}
\usepackage[utf8]{inputenc}
\usepackage[T1]{fontenc}
\usepackage{hyperref}
\title{ut-Hun Language Resources for K-12 Curriculum Development}
\author{Compiled by Claude Sonnet 5 (claude-sonnet-5) for Andr\'as Kornai}
\date{}
\begin{document}
\maketitle
\section{Introduction}
ut-Hun (ISO 639-3 \texttt{uth}) is a Kainji language of Nigeria. It is usually treated as the eastern dialect of Hun-Saare (Duka; Glottolog huns1239), and speakers have historically been called Dukkawa. The western dialect is Us-Saare (\texttt{uss}). Speaker estimates are about 50,000 in Rijau LGA, Niger State (a numerals source) and 73,000 (older Wikipedia figure, 1985). I found very little material.
\section{Primary Educational Resources}
I found no primers, textbooks or school programmes.
\section{Linguistic and Reference Materials}
\begin{itemize}
\item Hun-Saare language (Wikipedia): \url{https://en.wikipedia.org/wiki/Hun-Saare_language}
\item ut Hun, a numerals page: \url{https://lingweb.eva.mpg.de/channumerals/ut-Hun.htm}
\item ut hun, a Webonary page: \url{https://www.webonary.work/ut-hun/?p=9444}. I did not examine its contents, so its scope is unconfirmed.
\end{itemize}
\section{Religious and Community Resources}
\begin{itemize}
\item Scripture Earth entry for code uth: \url{https://www.scriptureearth.org/00eng.php?iso=uth}
\end{itemize}
\section{Digital and Audio Resources}
\begin{itemize}
\item Global Recordings Network page for ut-Hun: \url{https://globalrecordings.net/en/language/uth}
\end{itemize}
\section{Curriculum Development Context}
Only scripture and audio evangelistic materials and a few linguistic notes turned up. Any K-12 curriculum would need new work, ideally with the Hun community and with Dukkawa/Hun-Saare linguists.
\end{document}